\documentclass[11pt,a4paper]{article}

\usepackage[margin=2.6cm]{geometry}
\usepackage[T1]{fontenc}
\usepackage[utf8]{inputenc}
\usepackage{lmodern}
\usepackage{microtype}
\usepackage{amsmath,amssymb,mathtools}
\usepackage{booktabs}
\usepackage{enumitem}
\usepackage{hyperref}
\usepackage[round,authoryear]{natbib}

\title{A Two-Beam Theoretical Skeleton for Heterogeneous Learning Systems\\
\large Knowledge Organization and Competence Evolution}
\author{Technical research note}
\date{\today}

\newcommand{\Q}{\mathcal{Q}}
\newcommand{\A}{\mathcal{A}}
\newcommand{\E}{\mathcal{E}}

\begin{document}
\maketitle

\begin{abstract}
This technical note develops a minimal theoretical skeleton for studying systems in which
work allocation and competence evolution interact over time. It is not a paper and it does
not claim a new theory. Its purpose is to determine whether two independently established
theoretical beams can be combined without semantic or mathematical contradiction.

Beam~1 concerns the organization of distributed knowledge and the allocation or escalation
of problems, using Garicano's theory of knowledge hierarchies as the primary formal
foundation. Beam~2 concerns the evolution of competence through experience, learning,
retention and depreciation, drawing principally on learning-by-doing, task-specific human
capital, organizational learning, and Gutjahr's explicit competence-development model.

The central interface is operational experience: allocation determines which work is
performed; performed work generates experience; experience can change future competence;
and future competence changes the state on which subsequent allocation is based. The note
separates results established in the source literature from the synthesis introduced here,
and deliberately stops before developing a complete Heterogeneous Learning Systems model.
\end{abstract}

\section{Scope and question addressed}

The present question is deliberately narrower than the full Heterogeneous Learning Systems
(HLS) programme. We ask only whether the following two theoretical components can be
supported individually and joined coherently.

\begin{description}[leftmargin=2.8cm,style=nextline]
\item[Beam 1.] Given a distribution of problems and currently available knowledge or
competence, how should work and access to expertise be organized?
\item[Beam 2.] How can performing work change the competence available in the future?
\end{description}

The intended feedback structure is
\begin{equation}
\text{current state}
\longrightarrow
\text{work organization}
\longrightarrow
\text{operational experience}
\longrightarrow
\text{future competence}
\longrightarrow
\text{future work organization}.
\label{eq:cycle}
\end{equation}

The goal of this note is \emph{not} to solve the dynamic optimization problem implicit in
\eqref{eq:cycle}. It is to establish whether the arrows in this skeleton have defensible
theoretical foundations and whether the variables at their interfaces can be given compatible
meanings.

Three distinctions are important from the outset.

First, a \emph{knowledge set}---the set of problems for which a solution is available---is
not automatically the same object as a scalar or continuous \emph{competence score}.
Second, the cost of acquiring knowledge is not competence itself. Third, an allocation
decision and the experience produced by executing that allocation are related but distinct
objects.

These distinctions prevent a superficial joining of otherwise different source theories.

\section{Beam 1: organization of knowledge}

\subsection{Primary foundation}

The primary formal foundation is \citet{garicano2000}. Garicano studies production in which
workers encounter heterogeneous problems, knowledge acquisition is costly, and communication
allows a worker who does not know a solution to ask another worker.

Let
\[
Z\in\Q
\]
denote the problem encountered, with distribution \(F\) and density \(f\). A worker possesses
a knowledge set
\[
A\subseteq\Q,
\]
and solves the problem if \(Z\in A\).

In the autarkic case, Garicano orders problems from more to less frequent and considers a
knowledge interval \([0,Z_a]\). If the marginal cost of acquiring knowledge is \(c\), expected
net output is
\begin{equation}
E[y^a] = F(Z_a)-cZ_a,
\end{equation}
with interior first-order condition
\begin{equation}
f(Z_a)-c=0.
\end{equation}
Thus knowledge is acquired up to the point at which the marginal probability of solving an
additional problem equals its marginal acquisition cost.

\subsection{Communication and hierarchy}

Communication permits different workers to possess different knowledge sets. In Garicano's
model an organization assigns workers to classes, associates a knowledge set with each class,
and specifies whom a worker can ask when a solution is unknown. Asking for help consumes the
time of the worker being consulted. Denote this helping cost by \(h\).

The central trade-off is therefore between:
\begin{itemize}
\item acquiring knowledge locally, which is costly but avoids consultation; and
\item specializing knowledge and consulting other workers, which economizes on duplicated
knowledge acquisition but consumes communication and expert capacity.
\end{itemize}

Under the model assumptions, the resulting organization has a knowledge hierarchy:
production workers handle common problems and specialized problem solvers handle increasingly
exceptional problems. Information flows vertically and the organization exhibits management
by exception and a pyramidal structure.

This result is particularly relevant because the hierarchy is not imposed merely as an
organizational metaphor. It follows from the model's costs, problem frequencies and time
constraints.

\subsection{What Beam 1 establishes for the present skeleton}

Beam~1 supports the proposition that, for a given problem distribution and a given organization
of available knowledge, it can be inefficient for every operational unit to possess or use
all available expertise. Work can instead be organized so that common problems are handled
locally and exceptional problems are escalated.

In abstract notation, Beam~1 supplies a mapping of the form
\begin{equation}
a_t = \mathcal{R}(F_t,K_t,\kappa_t),
\label{eq:routing}
\end{equation}
where \(F_t\) describes problem demand, \(K_t\) represents the currently available organization
of knowledge, \(\kappa_t\) collects relevant costs and capacity constraints, and \(a_t\) is an
allocation/escalation organization.

Equation~\eqref{eq:routing} is an HLS-oriented abstraction, not an equation stated by
Garicano.

\subsection{What Beam 1 does not establish}

The base model explicitly assumes that the same problem recurs with probability zero and,
consequently, that communication does not itself generate learning. Garicano therefore does
not provide the temporal transition required for competence evolution.

Garicano also studies a knowledge set \(A\): whether solutions to classes of problems are
known. This should not be silently reinterpreted as a continuous measure of proficiency,
accuracy, speed or model quality.

Section IV of \citet{garicano2000} considers overlapping knowledge motivated by on-the-job
learning, but this does not constitute a dynamic law in which today's allocation endogenously
updates tomorrow's competence state. Beam~2 is required for that role.

\section{Beam 2: competence evolution through experience}

Beam~2 is not taken from a single reference. Its components are supported by complementary
literatures.

\subsection{Learning by doing}

The broad principle that productive activity can itself generate learning is classical.
\citet{arrow1962} formalized learning by doing as an economic mechanism through which
experience accumulated in production affects future productivity. For the present purpose,
the important proposition is modest: operation need not only consume current competence;
it can also contribute to future capability.

\subsection{Task specificity}

A generic stock of experience is too coarse for the intended interface with heterogeneous
problems. The task-specific human-capital literature, including \citet{gibbonswaldman2004},
supports the idea that experience accumulated in particular tasks can produce capability
that is especially relevant to those tasks.

This motivates indexing experience and competence by a task or problem region \(q\in\Q\),
without yet prescribing a particular representation of \(\Q\).

\subsection{Organizational learning: experience to knowledge}

\citet{argotemiron2011} provide the most direct conceptual bridge. They define organizational
learning around a change in organizational knowledge that occurs as a function of experience.
Their framework represents an ongoing cycle in which task-performance experience is converted
into knowledge; knowledge becomes embedded in organizational context; and the changed context
affects future experience.

The framework distinguishes three learning subprocesses:
\begin{enumerate}
\item knowledge creation from experience;
\item knowledge retention; and
\item knowledge transfer.
\end{enumerate}

For the present two-beam skeleton, creation and retention are already sufficient. Explicit
transfer can be added later and is not required to make the current feedback loop coherent.

Argote and Miron-Spektor also emphasize that experience should be characterized at a
fine-grained level. This is compatible with treating the experience produced by an allocation
as task-dependent rather than as a single undifferentiated count.

Their framework is conceptual rather than an optimization model. It does not provide a unique
state-transition equation for HLS.

\subsection{An explicit dynamic competence model}

As a primary current formal reference for this beam, \citet{gutjahr2011} provides a
quantitative model in which current work allocation changes future competence.

For project/task class \(i\) and period \(t\), let \(z_{it}\) be a competence score and
\(x_{it}\) the share of available work capacity invested in class \(i\). Gutjahr models
competence as
\begin{equation}
z_{it}
=
z_{i1}
-
\beta_i(t-1)
+
\eta_i\sum_{s=1}^{t-1}x_{is},
\label{eq:gutjahr}
\end{equation}
where \(\eta_i\) represents learning and \(\beta_i\) knowledge depreciation.

Competence is mapped to productive efficiency through a nondecreasing learning function
\begin{equation}
\gamma_{it}=\phi(z_{it}),\qquad 0\leq\gamma_{it}\leq 1.
\end{equation}
Gutjahr gives a concrete operational meaning to \(\gamma_{it}\): if a perfectly competent
team requires \(d\) units of real work time, a team with efficiency \(\gamma_{it}\) requires
\(d/\gamma_{it}\).

The essential dynamic mechanism is therefore
\begin{equation}
\text{work allocation}
\longrightarrow
\text{experience}
\longrightarrow
\text{competence}
\longrightarrow
\text{future productive efficiency}.
\end{equation}

The paper also demonstrates that the optimal pattern of concentration or mixing depends on
the learning function, depreciation and horizon. Thus learning does not mechanically imply
either specialization or diversification.

\subsection{What Beam 2 establishes for the present skeleton}

The combined literature supports the following ingredients:
\begin{enumerate}
\item task performance can generate learning;
\item learning can be task-specific;
\item knowledge/competence can be retained and can depreciate;
\item current work allocation can therefore alter future productive capability.
\end{enumerate}

A generic HLS-oriented transition can consequently be written as
\begin{equation}
C_{t+1}=\mathcal{L}(C_t,e_t;\theta_t),
\label{eq:learning}
\end{equation}
where \(e_t\) is learning-relevant exposure and \(\theta_t\) contains learning and retention
parameters. Real work and exposure coincide in Gutjahr's source model through its stated
work-dependent transition; an HLS translation must preserve or explicitly replace that
mapping. Equation~\eqref{eq:learning} is a synthesis, not a formula asserted by the cited
literature.

\subsection{What Beam 2 does not establish}

Gutjahr's competence is aggregated at team level in the basic model. It is not a theory of
routing among heterogeneous machine-learning models. Nor does the cited foundation determine
how knowledge should be transferred from a teacher model to a learner, how machine-learning
competence should be measured, or how experience generalizes across neighboring task regions.

Those questions lie outside the present note.

Gutjahr is not the theory of Beam~2. Its formalization is one source among multiple
independent lines, and its team-level state, work variable, and efficiency mapping must not
be imposed on HLS without an explicit ontology mapping.

\section{Compatibility of the two beams}

The two beams address different questions and use different state semantics. This is an
advantage if the interface is chosen carefully.

\subsection{Objects that should not be identified}

Three tempting identifications should be rejected.

\paragraph{Knowledge-acquisition cost is not competence.}
Garicano's \(c\) is the marginal cost of acquiring knowledge. Gutjahr's \(z_{it}\) is a
competence score and \(\phi(z_{it})\) is productive efficiency. A relation such as
\(c_t=\psi(C_t)\) is not supplied by either theory and is unnecessary for joining them.

\paragraph{Knowledge set is not automatically proficiency.}
Garicano's
\[
A_t\subseteq\Q
\]
records which problem solutions are known. A continuous competence variable \(C_t(q)\)
can instead represent how effectively a known or learnable class of tasks is performed.
Defining
\[
A_t=\{q:C_t(q)\geq\vartheta\}
\]
would be an additional modeling assumption, not a consequence of the foundations reviewed
here.

\paragraph{Allocation, real work and learning exposure are distinct.}
An allocation decision determines what work is attempted and by whom. Real work is the
resource actually consumed in execution. Learning exposure is the part and type of the
resulting experience relevant to a competence transition. Keeping these objects distinct
allows failures, partial completion, feedback quality, support work, or non-learning work
without changing the conceptual skeleton.

\subsection{Compatible common structure}

Beam~1 ends with an operational allocation:
\[
(F_t,K_t,\kappa_t)\longrightarrow a_t.
\]
Execution maps allocation to real work:
\[
w_t=\mathcal{W}(a_t,F_t,\text{execution conditions}).
\]
Work and outcomes map to learning-relevant exposure:
\[
e_t=\mathcal{E}(w_t,\text{execution outcomes}),
\]
and Beam~2 uses that exposure in a competence transition:
\[
C_{t+1}=\mathcal{L}(C_t,e_t).
\]

The maps \(\mathcal W\), \(\mathcal E\), and \(\mathcal L\) are part of an HLS
instantiation; they are not supplied jointly by the source papers. The interface is therefore
a candidate modeling scaffold governed by the canonical HLS ontology, not an established
cross-paper identification.

\section{The Beam 1--Beam 2 interface}

The smallest defensible interface is
\begin{equation}
(F_t,S_t)
\xrightarrow{\;\mathcal R\;}
a_t
\xrightarrow{\;\mathcal W\;}
w_t
\xrightarrow{\;\mathcal E\;}
e_t
\xrightarrow{\;\mathcal L\;}
S_{t+1},
\label{eq:interface}
\end{equation}
where \(S_t\) is an intentionally generic state containing whatever current knowledge and
competence variables are required.

The use of \(S_t\), rather than prematurely collapsing everything into one scalar competence,
is important. A future model may require at least two components:
\begin{equation}
S_t=(K_t,C_t),
\end{equation}
where \(K_t\) represents the availability of problem-solution knowledge and \(C_t\) represents
proficiency or efficiency. The current theoretical foundations do not yet justify eliminating
either component.

Here \(w_t\) is real operational work and \(e_t\) is learning-relevant exposure. They may
coincide in a source model or special HLS instantiation, but they are not identified by
default. The maps \(\mathcal W\), \(\mathcal E\), and \(\mathcal L\) must state their
types, units, domains, and assumptions.

\subsection{A task-indexed learning abstraction}

If attention is restricted to the intensive competence margin, a simple representation
consistent with the mechanisms of Beam~2 is
\begin{align}
E_{t+1}(q) &= \rho_q E_t(q) + e_t(q),\\
C_t(q) &= L_q(E_t(q)),
\label{eq:experience}
\end{align}
where \(E_t(q)\) is retained accumulated exposure, \(e_t(q)\) is current
learning-relevant exposure, \(\rho_q\) is a retention factor and \(L_q\) is a learning
function.

This is \textbf{not} presented as an established equation from Arrow, Argote, Gibbons--Waldman
or Gutjahr. It is a minimal synthesis that preserves their relevant mechanisms: experience,
task specificity, learning and depreciation.

\subsection{Source support and remaining mapping burden}

The interface in \eqref{eq:interface} requires only two substantive claims.

First, operational organization determines which tasks are performed or escalated. This is
the subject of Beam~1.

Second, task-performance experience can alter future knowledge or competence. This is
supported independently by Beam~2 and is made particularly explicit by the experience-to-
knowledge cycle of \citet{argotemiron2011} and the dynamic competence equation of
\citet{gutjahr2011}.

These source results make the scaffold plausible, but they do not establish the particular
maps or their semantic equivalence in HLS. The bridge remains \textbf{CANDIDATE /
HYPOTHESIS}; integration must proceed through the canonical ontology rather than direct
paper-to-paper variable identification.

\section{Minimal combined skeleton}

We can now write the partial combination without claiming a complete HLS theory.

Let \(F_t\) denote the distribution of incoming problems and \(S_t\) the current
knowledge/competence state. The operational organization is
\begin{equation}
a_t=\mathcal R(F_t,S_t,\kappa_t).
\end{equation}
Executing this organization generates experience
\begin{equation}
x_t=\mathcal X(a_t,F_t,\omega_t),
\end{equation}
where \(\omega_t\) denotes execution outcomes or other factors determining realized
experience. Competence/knowledge then evolves according to
\begin{equation}
S_{t+1}=\mathcal L(S_t,x_t,\theta_t).
\end{equation}
The next period therefore uses
\begin{equation}
a_{t+1}=\mathcal R(F_{t+1},S_{t+1},\kappa_{t+1}).
\end{equation}

Combining the equations gives the feedback system
\begin{equation}
S_{t+1}
=
\mathcal L\!\left(
S_t,
\mathcal X\!\left(
\mathcal R(F_t,S_t,\kappa_t),
F_t,\omega_t
\right),
\theta_t
\right).
\label{eq:combined}
\end{equation}

Equation~\eqref{eq:combined} is the central mathematical statement of this note. It should be
read as a structural composition, not as a solved model. Each operator still requires a
specific instantiation before optimization or theorem development is possible.

\subsection{What the combination adds}

Neither beam alone contains the full feedback.

Beam~1 can organize current knowledge efficiently but, in its base form, does not let the
organization's own operation change future competence.

Beam~2 allows current allocation to change future competence but does not supply Garicano's
problem-matching and escalation structure.

Their combination introduces the possibility that an operationally sensible allocation at
time \(t\) changes the competence state used to organize work at \(t+1\). This is the precise
sense in which present organization and future capability become linked.

No stronger claim is needed at this stage.

\section{Epistemic status of the construction}

To avoid attributing our synthesis to the source literature, Table~\ref{tab:status} separates
established results from the present construction.

\begin{table}[ht]
\centering
\small
\begin{tabular}{p{0.30\linewidth}p{0.20\linewidth}p{0.39\linewidth}}
\toprule
Statement & Status & Basis\\
\midrule
Costly knowledge acquisition and communication can induce knowledge hierarchies and
management by exception.
& Established
& Garicano (2000).\\[2mm]

Productive experience can generate learning.
& Established
& Learning-by-doing literature; Arrow (1962).\\[2mm]

Learning can be task-specific.
& Established
& Task-specific human-capital literature; Gibbons and Waldman (2004).\\[2mm]

Organizational knowledge can change as a function of task-performance experience; knowledge
is created, retained and transferred.
& Established framework
& Argote and Miron-Spektor (2011).\\[2mm]

Work allocation can explicitly affect future competence through learning and depreciation.
& Established formal model
& Gutjahr (2011), under its assumptions.\\[2mm]

Allocation can determine real work, which may generate task-dependent learning exposure
entering a competence transition.
& Candidate / hypothesis
& Requires explicit \(\mathcal W\), \(\mathcal E\), and \(\mathcal L\) mappings in the HLS ontology.\\[2mm]

\((F_t,S_t)\to a_t\to w_t\to e_t\to S_{t+1}\to a_{t+1}\).
& Candidate HLS scaffold
& Not stated as such by any single cited source.\\[2mm]

A particular machine-learning implementation of \(S_t,\mathcal R,\mathcal W,\mathcal E,\mathcal L\).
& Open
& Requires later modeling and empirical work.\\
\bottomrule
\end{tabular}
\caption{Epistemic status of the two-beam skeleton.}
\label{tab:status}
\end{table}

The individual ingredients are source-grounded. Their mapping and combination should be
described as a \emph{candidate modeling scaffold}, not as an established theorem or a
completed transfer inherited from the literature.

\section{Consistency checks}

A useful skeleton should degrade sensibly when one of its mechanisms is removed.

\subsection{No competence evolution}

If
\[
\mathcal L(S_t,x_t,\theta_t)=S_t,
\]
operation does not change competence. The feedback through learning disappears and the
problem reduces to repeated organization of a fixed knowledge/competence state. This is the
natural static limit of Beam~1.

\subsection{No effect of allocation on experience}

If realized experience \(x_t\) is independent of \(a_t\), then current work organization does
not influence future competence through the proposed interface. The two beams coexist but are
not dynamically coupled. This is an important null case.

\subsection{No depreciation}

Setting retention to one in a representation such as
\[
e_{t+1}(q)=\rho_q e_t(q)+x_t(q)
\]
removes forgetting while retaining learning by experience. This corresponds to a standard
limiting case of competence-development models.

\subsection{No learning}

Setting the learning contribution to zero eliminates the Beam~2 feedback. A dynamic controller
should then have no competence-development reason to sacrifice current operational value.

\subsection{Why these checks matter}

These reductions show that the combined skeleton does not require learning to manufacture an
effect. If the relevant dynamic mechanisms are absent, the dynamic coupling disappears. This
is a basic but necessary safeguard against an artificial construction.

\section{Current limits and stopping point}

The present armazón is intentionally incomplete.

It does not yet specify:
\begin{itemize}
\item how machine-learning competence should be represented or measured;
\item whether knowledge availability \(K_t\) and proficiency \(C_t\) can be reduced to a
single state without loss;
\item how experience on one task generalizes to related tasks;
\item how explicit knowledge transfer between models should operate;
\item how several heterogeneous learning agents should be jointly managed;
\item how costs, latency and capacity should enter a complete HLS objective;
\item how non-stationary problem demand should be modeled;
\item or what optimal dynamic policy follows from the combined system.
\end{itemize}

These omissions do not invalidate the partial combination. They delimit it.

At the present stage the defensible conclusion is narrower:

\begin{quote}
A coherent theoretical skeleton exists in which the organization of current knowledge
determines operational allocation, operational allocation generates task-performance
experience, and experience can modify the competence state on which future organization
depends. The individual mechanisms are grounded in established theories, while their
composition into this feedback loop is the present synthesis.
\end{quote}

This is the appropriate stopping point before extending the construction toward the full HLS
problem.


\bibliographystyle{plainnat}
\bibliography{references}

\end{document}
